\section{Metodología}
La metodología implementada para la realización del sistema, se fundamentó con el 
marco de trabajo Scrum, reconocido por su flexibilidad y capacidad de adaptación [6]. Se 
organizaron ciclos de trabajo denominados sprints de dos semanas, en los cuales el 
equipo debía entregar incrementos funcionales del sistema. Se realizaron reuniones de 
planificación, reuniones diarias (daily scrum) para el seguimiento del progreso, revisiones 
de sprint y retroalimentación, con el fin de garantizar una mejora continua. Este enfoque 
se complementa con documentación formal y un modelo de especificación de 
requerimientos detallado, tal como se establece en el Informe de Requerimientos [10]. 

Desde el punto de vista técnico, la elección de las herramientas y los frameworks se 
realizó con base en análisis y comparaciones realizadas. Laravel fue adoptado como 
framework backend debido a su simplicidad, seguridad integrada y ecosistema de 
herramientas como Eloquent ORM, Blade y migraciones [1]. Aunque Symfony se 
consideró por su robustez y orientación empresarial [7], se determinó que Laravel 
respondía mejor a los objetivos del proyecto. En el frontend, React fue implementado por 
su capacidad de reutilización de componentes, eficiencia y soporte por parte de una 
amplia comunidad [2]. Estudios comparativos demuestran que React ofrece un balance 
adecuado entre curva de aprendizaje y flexibilidad, en contraste con Angular, que 
presenta mayor complejidad [4], [5]. 

La aplicación móvil se diseñó únicamente para el sistema operativo Android, tomando en 
cuenta la accesibilidad de esta plataforma en el contexto del proyecto. Se contempló la 
futura incorporación de React Native para extender la compatibilidad hacia iOS y 
aprovechar las ventajas de compartir código entre plataformas [3]. 

El sistema está constituido por tres módulos fundamentales: la base de datos centralizada, 
la aplicación móvil y el portal web. El portal web facilita la administración, el registro de 
tickets y la creación de informes; la aplicación móvil brinda a los técnicos la posibilidad 
de recibir notificaciones push y guardar evidencias multimedia; y contar con una base de 
datos centralizada garantiza que la información sea íntegra, se pueda rastrear y esté 
disponible. La integración de estos módulos se llevó a cabo de acuerdo con el modelo 
arquitectónico MVC [8], lo cual permitió dividir las responsabilidades y hacer que el 
sistema pueda escalar a un  futuro. 